The Hubble ``tension'' has persisted for more than a decade and significantly shifted focus amongst endeavors to precisely measure \Hcst. Before $\sim$2012, the primary motivation for measuring $H_0$ with increasing precision was to provide a strong prior for constraining $\Lambda$CDM parameters and understanding dark energy \citep[e.g.,][]{Suyu:2012}. As the statistical significance of the discrepancy between early- and late-Universe measurements grew, relentless attention was given to systematics of local distance measurements with an implicit expectation that the tension would be resolved once the responsible systematic error was identified. In other words, there was a non-negligible preference among many to discard direct measurements (even ones with small and detailed uncertainties) based on astronomical techniques in favor of the \Hcst\ value derived from the perceived ``simple physics'' of the early Universe and a specific cosmological model. A narrative of merely two or three direct measurements disagreeing with one extremely precise prediction based on the {\it Planck} CMB and flat $\Lambda$CDM lent credence to this view, despite a plethora of semi-(in)dependent measurements painting similar pictures \citep[e.g.,][and references therein]{CosmoVerse:2025txj}. A statistically rigorous combination of the available measurements was thus needed to resolve whether any one method might be deemed `bad' or unreliable. The present work provides this long overdue combination supported by expert knowledge from all relevant disciplines, and the results presented discredit the expectation of `unknown unknowns' capable of misleading the direct measurement of the local expansion rate using astronomical techniques.

The {\DN} framework introduced here establishes a reproducible and extensible methodology for synthesizing local distance constraints. Its transparency and modularity allow for the straightforward integration of new data, methods, or calibrators, enabling the community to track the evolution of $H_0$ estimates over time. In addition to the primary distance ladder, the network already incorporates several fully independent paths to $H_0$, including the maser-based anchor, the TRGB–SBF route, and SNe~II distances based on astrophysical modeling. These methods rely on distinct physical assumptions and observational systematics, yet yield results consistent with the primary solution, reinforcing the robustness of the empirical determination. While current uncertainties for these paths are larger, their future competitiveness will depend on targeted improvements in calibration, modeling, and sample size. Promising opportunities include searches for SNe~II in SBF or SNe~Ia host galaxies, extending TRGB and Cepheid measurements deeper into the Hubble flow using ELTs, denser tiling of NGC~4258 to refine JAGB and TRGB calibrations, and establishing additional anchors via geometric methods \citep[e.g., using population-II pulsating stars calibrated by parallax, e.g.,][]{2025arXiv250916331L}. Strengthening these independent routes and developing new anchors will be critical for validating and advancing precision cosmology. The approach presented here serves both as a benchmark of current capabilities and as a foundation for future consensus efforts in the determination of $H_0$.

The {\DN} presented in this work provides the most precise direct measurement to date of the Hubble constant, $H_0 = \Hvaluebase \pm \Herrorbase \, \Hunit$, achieving a relative uncertainty of {\Hrelerrorbase}. This value is the result of a comprehensive, covariance-weighted combination of local distance measurements incorporating multiple distance indicators, each rigorously assessed for reliability, consistency, and methodological transparency. By explicitly accounting for shared uncertainties and correlations, the {\DN} formalism surpasses the precision of previous individual or pairwise combinations and reflects a collaborative consensus on data selection, treatment, and analysis strategy.

In constructing the distance network, certain methodological deviations from state-of-the-art analyses were necessary to achieve full covariance treatment across indicators. For example, the use of supernova intercepts for calibration simplifies network integration but does not incorporate the full cross-rung covariance present in datasets such as Pantheon+, as done in \citet{2022ApJ...934L...7R}. These tradeoffs are explicitly documented to ensure transparency and reproducibility. We are publicly releasing the full {\DN} code upon acceptance of this paper at \url{https://github.com/StefCas789/H0DN}. We are planning to additionally publish the code in the Astrophysics Source Code Library\footnote{\url{https://ascl.net}} -- the corresponding DOI will be updated and linked here in a future version of the paper.

A central feature of our approach is the construction of a baseline solution exclusively based on top-ranked calibrators and tracers, as determined by agreement among all authors as representatives of their respective communities. This was complemented by a suite of predefined variants, decided prior to conducting the analysis, which explore the sensitivity of the result to alternative assumptions, datasets, or methodological choices. The resulting spread in $H_0$ values across all variants is significantly smaller than the statistical uncertainty of any individual variant, confirming that no single dataset or technique dominates the result and that the network solution is not driven by any particular methodology or systematic effect. This is also confirmed by investigating routes through the network that are designed to not have any elements in common, leading to results that are highly consistent with the baseline analysis.

Further scrutiny of systematics affecting {\DN} elements is both necessary and useful for further improvements to {\Hcst} accuracy. However, reconciling the {\DN}-based direct {\Hcst} measurement with the expectation from $\Lambda$CDM would require an extremely unlikely alignment of systematics conspiring to overestimate {\Hcst}. The diversified portfolio of methods and datasets incorporated into the {\DN} effectively protects against significant bias and demonstrates that the Hubble tension cannot be eliminated by invoking as-yet unidentified systematics in any single method or dataset.

Our result differs by more than $6\sigma$ from the value predicted by the $\Lambda$CDM model calibrated to \textit{Planck} CMB observations, by more than $7\sigma$ from that calibrated to SPT+ACT CMB measurements, which are given as  $\Hcst = \Hvaluecmb \pm \Herrorcmb \Hunit $ from Planck+SPT+ACT \citep[Eq.~(54)~of][]{2025arXiv250620707C}, and by about $5\sigma$ from the combination of BBN and DESI BAO constraints, which are $\Hcst = 68.51 \pm 0.58 \Hunit$   \citep[Tab.~V~of][]{DESI:2025zgx}.

Since the {\DN} provides a conclusive direct \Hcst\ measurement with demonstrated control of systematics, the resulting significant discrepancy retires the aforementioned expectation that identifying a simple systematic could resolve the Hubble constant tension. 
In this vein, we advocate that the empirically determined value of $H_0$ from the {\DN} be used as a prior in cosmological analyses. 
This marks a significant shift in perspective: rather than serving solely to constrain dark energy models, as envisioned a decade ago, the improved accuracy of $H_0$ now exposes a broader inconsistency within the standard cosmological framework and strengthens the case for new physics or a deeper reassessment of early-Universe inferences. The evolving role of $H_0$ has already reshaped our understanding of precision cosmology, and further surprises may lie ahead.
